\chapter{Riemannian Geometry Background}

\section{Notation and conventions}

Throughout this chapter, $X$, $Y$, $W$, and $Z$ are fixed vector fields on a manifold $\M$, and $A$ and $B$ are tensor fields, possibly of different types. We assume that $\M$ carries a Riemannian metric $g$, or a smooth family of such metrics depending on a parameter $t$. The metric extends to arbitrary tensors, and $\nabla$ denotes the Levi-Civita connection of $g$.

The connection is metric compatible and satisfies the product rule
\[
\nabla_X(A \otimes B) = (\nabla_X A) \otimes B + A \otimes (\nabla_X B).
\]
We write $\nabla_X A$ as $\nabla A(X,\ldots)$ when $\nabla$ is used without a subscript, so that $X$ is the first slot.

\begin{definition}[\dcref{2.1.1}]
\label{topping:notn-second-covariant-derivative}
\group{topping-background}
\source{topping:page-0026}
The second covariant derivative of a tensor field is
\[
\nabla^2_{X,Y} := \nabla_X \nabla_Y - \nabla_{\nabla_X Y}.
\]
\end{definition}

\begin{definition}[\dcref{2.1.1}]
\label{topping:notn-hessian}
\group{topping-background}
\source{topping:page-0026}
For a function $f : \M \to \mathbb{R}$, the Hessian is
\[
\Hess(f) := \nabla df.
\]
\end{definition}

\begin{definition}[\dcref{2.1.1}]
\label{topping:notn-rough-laplacian}
\group{topping-background}
\source{topping:page-0026}
The connection Laplacian of a tensor field $A$ is
\[
\Delta A := \tr_{12} \nabla^2 A.
\]
\end{definition}

\begin{definition}[\dcref{2.1.1}]
\label{topping:notn-curvature-operator}
\group{topping-background}
\source{topping:page-0026}
The curvature tensor is defined by
\[
R(X,Y) := \nabla_Y \nabla_X - \nabla_X \nabla_Y + \nabla_{[X,Y]}.
\]
\end{definition}

\begin{definition}[\dcref{2.1.1}]
\label{topping:notn-riemann-tensor}
\group{topping-background}
\source{topping:page-0026}
The Riemann curvature tensor is
\[
\Rm(X,Y,W,Z) := \langle R(X,Y)W, Z \rangle.
\]
\end{definition}

\begin{definition}[\dcref{2.1.1}]
\label{topping:notn-richi-tensor}
\group{topping-background}
\source{topping:page-0027}
The Ricci tensor and scalar curvature are
\[
\Ric(X,Y) := \tr \Rm(X,\cdot,Y,\cdot), \qquad R := \tr \Ric.
\]
\end{definition}

\begin{definition}[\dcref{2.1.1}]
\label{topping:notn-star-product}
\group{topping-background}
\source{topping:page-0027}
The notation $A * B$ denotes any tensor field obtained as a universal linear combination of tensors built from $A \otimes B$ using metric contractions, index raising and lowering, and permutations of tensor factors.
\end{definition}

\begin{definition}[\dcref{2.1.1}]
\label{topping:notn-divergence}
\group{topping-background}
\source{topping:page-0028}
The divergence of a tensor field $T \in \Gamma(\otimes^k T^*\M)$ is
\[
\delta(T) := - \tr_{12} \nabla T.
\]
\end{definition}

\begin{definition}[\dcref{2.1.1}]
\label{topping:notn-gravitation-tensor}
\group{topping-background}
\source{topping:page-0028}
For a symmetric $2$-tensor $T$, the gravitation tensor is
\[
G(T) := T - \frac12 (\tr T) g.
\]
\end{definition}

\begin{remark}[\dcref{2.1.1}]
\label{topping:notn-background-remark}
\group{topping-background}
\source{topping:page-0026,topping:page-0027,topping:page-0028}
The exposition fixes the sign conventions
\[
\nabla^2_{X,Y} = \nabla_X \nabla_Y - \nabla_{\nabla_X Y}, \qquad
R(X,Y) = \nabla_Y \nabla_X - \nabla_X \nabla_Y + \nabla_{[X,Y]},
\]
and the shorthand identities
\[
R(.,.)A = A * \Rm, \qquad \nabla(A * B) = (\nabla A) * B + A * (\nabla B).
\]
The contracted Bianchi identity is written as
\[
\delta G(\Ric) = \delta \Ric + \frac12 dR = 0.
\]
\end{remark}

\section{Einstein metrics}

\begin{proposition}[\dcref{2.2}]
\label{topping:einstein-metrics-const-scalar}
\group{topping-background}
\source{topping:page-0029}
If $\Ric(g) = \lambda g$ for a function $\lambda : \M \to \mathbb{R}$, then in dimensions $n \neq 2$ the scalar curvature is constant and $\lambda$ is constant, so the two Einstein conventions agree.
\end{proposition}

\section{Deformation of geometric quantities}

Suppose $g = g_t$ is a smooth family of metrics and $h := \partial g_t / \partial t$.

\begin{proposition}[\dcref{2.3.1}]
\label{topping:variation-levi-civita}
\group{topping-background}
\source{topping:page-0029}
For the variation tensor $\Pi(X,Y) := \partial (\nabla_X Y)/\partial t$, one has
\[
\langle \Pi(X,Y), Z \rangle
= \frac12 \bigl[(\nabla_Y h)(X,Z) + (\nabla_X h)(Y,Z) - (\nabla_Z h)(X,Y)\bigr].
\]
\end{proposition}

\begin{remark}[\dcref{2.3.2}]
\label{topping:variation-vector-field}
\group{topping-background}
\source{topping:page-0030}
If $V = V(t)$ is a $t$-dependent vector field, then
\[
\frac{\partial}{\partial t}\nabla_X V = \Pi(X,V) + \nabla_X \frac{\partial V}{\partial t}.
\]
\end{remark}

\begin{remark}[\dcref{2.3.3}]
\label{topping:variation-covariant-derivative-star}
\group{topping-background}
\source{topping:page-0030}
For a fixed tensor field $A$, the variation of the covariant derivative satisfies
\[
\frac{\partial}{\partial t}\nabla A = A * \nabla h,
\]
and for a $t$-dependent tensor field,
\[
\frac{\partial}{\partial t}\nabla A - \nabla \frac{\partial A}{\partial t} = A * \nabla h.
\]
\end{remark}

\begin{proposition}[\dcref{2.3.4}]
\label{topping:variation-curvature-operator}
\group{topping-background}
\source{topping:page-0030}
The curvature operator varies by
\[
\frac{\partial}{\partial t} R(X,Y)W = (\nabla_Y \Pi)(X,W) - (\nabla_X \Pi)(Y,W).
\]
\end{proposition}

\begin{proposition}[\dcref{2.3.5}]
\label{topping:variation-riemann-tensor}
\group{topping-background}
\source{topping:page-0031}
The Riemann curvature tensor varies according to
\[
\frac{\partial}{\partial t}\Rm(X,Y,W,Z)
= \frac12\bigl[h(R(X,Y)W,Z) - h(R(X,Y)Z,W)\bigr]
\]
\[
\qquad\qquad
 + \frac12\bigl[\nabla^2_{Y,W} h(X,Z) - \nabla^2_{X,W} h(Y,Z)
 + \nabla^2_{X,Z} h(Y,W) - \nabla^2_{Y,Z} h(X,W)\bigr].
\]
\end{proposition}

\begin{proposition}[\dcref{2.3.6}]
\label{topping:variation-trace}
\group{topping-background}
\source{topping:page-0031}
For any $t$-dependent tensor $\alpha \in \Gamma(\otimes^2 T^*\M)$,
\[
\frac{\partial}{\partial t}(\tr \alpha) = -\langle h,\alpha\rangle + \tr\frac{\partial \alpha}{\partial t}.
\]
\end{proposition}

\begin{definition}[\dcref{2.3.7}]
\label{topping:variation-lichnerowicz}
\group{topping-background}
\source{topping:page-0031}
The Lichnerowicz Laplacian on symmetric $2$-tensors is
\[
(\Delta_{\mathcal L} h)(X,W) := (\Delta h)(X,W) - h(X,\Ric(W)) - h(W,\Ric(X)) + 2\tr h(R(X,\cdot)W,\cdot).
\]
\end{definition}

\begin{proposition}[\dcref{2.3.7}]
\label{topping:variation-ricci}
\group{topping-background}
\source{topping:page-0031,topping:page-0032}
The Ricci tensor varies by
\[
\frac{\partial}{\partial t}\Ric
= -\frac12 \Delta_{\mathcal L} h - \frac12 \mathcal L_{(\delta G(h))^\#} g.
\]
Equivalently,
\[
\frac{\partial}{\partial t}\Ric(X,W)
= \Delta \Ric(X,W) - 2\langle \Ric(X),\Ric(W)\rangle + 2\langle \Rm(X,\cdot,W,\cdot), \Ric\rangle.
\]
\end{proposition}

\begin{remark}[\dcref{2.3.8}]
\label{topping:lie-derivative-symmetric-gradient}
\group{topping-background}
\source{topping:page-0032}
For a $1$-form $\omega$,
\[
\mathcal L_{\omega^\#} g(X,W) = \nabla\omega(X,W) + \nabla\omega(W,X).
\]
If $\omega = df$, then
\[
\mathcal L_{(df)^\#} g = 2\Hess(f).
\]
\end{remark}

\begin{proposition}[\dcref{2.3.9}]
\label{topping:variation-scalar-curvature}
\group{topping-background}
\source{topping:page-0032,topping:page-0033,topping:page-0037}
The scalar curvature varies by
\[
\frac{\partial R}{\partial t} = -\langle h,\Ric\rangle + \delta^2 h - \Delta(\tr h).
\]
\end{proposition}

\begin{proposition}[\dcref{2.3.10}]
\label{topping:variation-divergence-one-form}
\group{topping-background}
\source{topping:page-0033,topping:page-0038}
For any $t$-dependent $1$-form $\omega$,
\[
\frac{\partial}{\partial t}\delta\omega
= \delta\frac{\partial \omega}{\partial t} + \langle h,\nabla\omega\rangle - \langle \delta G(h),\omega\rangle.
\]
\end{proposition}

\begin{proposition}[\dcref{2.3.11}]
\label{topping:variation-divergence-gravitation}
\group{topping-background}
\source{topping:page-0033,topping:page-0039,topping:page-0040}
If $T \in \Gamma(\Sym^2 T^*\M)$ is independent of $t$, then
\[
\left(\frac{\partial}{\partial t}\delta G(T)\right) Z
= -T\bigl((\delta G(h))^\#, Z\bigr) + \langle h, \nabla T(\cdot,\cdot; Z) - \frac12 \nabla_Z T\rangle.
\]
\end{proposition}

\begin{proposition}[\dcref{2.3.12}]
\label{topping:variation-volume-form}
\group{topping-background}
\source{topping:page-0033,topping:page-0040}
The volume form varies by
\[
\frac{\partial}{\partial t} dV = \frac12 (\tr h)\, dV.
\]
\end{proposition}

\section{Laplacian of the curvature tensor}

\begin{proposition}[\dcref{2.4.1}]
\label{topping:laplacian-curvature-tensor}
\group{topping-background}
\source{topping:page-0041,topping:page-0042}
The rough Laplacian of the curvature tensor satisfies
\[
(\Delta \Rm)(X,Y,W,Z)
= -\nabla^2_{Y,W}\Ric(X,Z) + \nabla^2_{X,W}\Ric(Y,Z)
 - \nabla^2_{X,Z}\Ric(Y,W) + \nabla^2_{Y,Z}\Ric(X,W)
\]
\[
\qquad\qquad
- \Ric(R(W,Z)Y,X) + \Ric(R(W,Z)X,Y)
- 2\bigl(B(X,Y,W,Z) - B(X,Y,Z,W)
 + B(X,W,Y,Z) - B(X,Z,Y,W)\bigr).
\]
\end{proposition}

\section{Evolution under Ricci flow}

\begin{proposition}[\dcref{2.5.1}]
\label{topping:ricci-flow-curvature-evolution}
\group{topping-background}
\source{topping:page-0042,topping:page-0043}
Under the Ricci flow,
\[
\frac{\partial}{\partial t}\Rm = \Delta \Rm + \Rm * \Rm.
\]
More explicitly,
\[
\frac{\partial}{\partial t}\Rm(X,Y,W,Z)
= (\Delta \Rm)(X,Y,W,Z)
- \Ric(R(X,Y)W,Z) + \Ric(R(X,Y)Z,W)
\]
\[
\qquad\qquad
- \Ric(R(W,Z)X,Y) + \Ric(R(W,Z)Y,X)
+ 2\bigl(B(X,Y,W,Z) - B(X,Y,Z,W)
+ B(X,W,Y,Z) - B(X,Z,Y,W)\bigr).
\]
\end{proposition}

\begin{remark}[\dcref{2.5.2}]
\label{topping:ricci-flow-curvature-heat}
\group{topping-background}
\source{topping:page-0043}
The curvature tensor evolves by a heat-type equation,
\[
\frac{\partial}{\partial t}\Rm = \Delta \Rm + \Rm * \Rm.
\]
\end{remark}

\begin{proposition}[\dcref{2.5.3}]
\label{topping:ricci-flow-ricci-evolution}
\group{topping-background}
\source{topping:page-0043}
Under the Ricci flow, the Ricci tensor evolves by
\[
\frac{\partial}{\partial t}\Ric = \Delta_{\mathcal L}(\Ric),
\]
equivalently
\[
\frac{\partial}{\partial t}\Ric(X,W)
= \Delta \Ric(X,W) - 2\langle \Ric(X), \Ric(W)\rangle + 2\langle \Rm(X,\cdot,W,\cdot), \Ric\rangle.
\]
\end{proposition}

\begin{proposition}[\dcref{2.5.4}]
\label{topping:ricci-flow-scalar-evolution}
\group{topping-background}
\source{topping:page-0043}
Under the Ricci flow, the scalar curvature evolves by
\[
\frac{\partial R}{\partial t} = \Delta R + 2|\Ric|^2.
\]
\end{proposition}

\begin{remark}[\dcref{2.5.4}]
\label{topping:ricci-flow-scalar-ineq}
\group{topping-background}
\source{topping:page-0043}
Writing $\Ric = \mathring{\Ric} + \frac{R}{n}g$, one gets the pointwise inequality
\[
|\Ric|^2 \ge \frac{R^2}{n}.
\]
\end{remark}
